\section{結言}
この観測窓では、新規モデルより「AIを実際のシステムとして運用したときに何が起きるか」が中心テーマになった。評価エージェントの隔離、neocloud、infostealer、利用枠やモデル供給の変更など、能力競争と同時に運用境界の設計が重要になっている。

一方、Redwood、PuRo-2B、CLEARは設計自動化、低コスト学習、安全性更新の効率化を示した。M\&Aやインフラ調達も含め競争軸は計算基盤、配布層、セキュリティ、ロボティクスへ広がっている。未確認の買収・価格・性能値は、X上で共有された主張として読む必要がある。
